\SourceLeaf{300}
\par\bigskip\pdfbookmark[0]{De nexu. Fragmentum X}{nexus-fragment-X}\centerline{[X.]}\medskip
\begin{center}\emph{Vertheilung der quadratischen Reste in Zwölftel.}\end{center}
$p$ Primzahl; $(r)$ Anzahl der quadratischen Reste von $p$, welche zwischen $\frac{r-1}{12}p$ und $\frac r{12}p$ liegen.

Erster Fall; $p=24n+1$.
\[
\begin{aligned}
2t&\quad\text{Anzahl der Classen für den Determinans }-p\\
4u&\quad\text{Anzahl der Classen für den Determinans }-3p
\end{aligned}
\]
\[
\begin{aligned}
(1)=(12)&=\tfrac16(6n+3t+2u)\\
(2)=(4)=(6)=(7)=(9)=(11)&=\tfrac16(6n-3t+2u)\\
(3)=(5)=(8)=(10)&=\tfrac16(6n+3t-4u)
\end{aligned}
\]
\[
\begin{array}{r|r|r|r|r|r|r}
p&n&t&u&(1)&(2)&(3)\\\hline
73&3&2&3&5&3&2\\
97&4&2&3&6&4&3\\
193&8&2&6&11&9&5\\
241&10&6&3&14&8&11
\end{array}
\]

Zweiter Fall; $p=24n+13$.
\[
\begin{aligned}
2t&\quad\text{Anzahl der Classen für den Determinans }-p;\\
4u&\quad\text{Anzahl der Classen für den Determinans }-3p.
\end{aligned}
\]
\[
\begin{aligned}
(1)=(3)=(10)=(12)&=\tfrac12(2n+1+t)\\
(2)=(6)=(7)=(1.1)&=\tfrac12(2n+1-t)\\
(4)=(9)&=\tfrac12(2n+1-t+2u)\\
(5)=(8)&=\tfrac12(2n+1+t-2u)
\end{aligned}
\]
\[
\begin{array}{r|r|r|r|r|r|r|r}
p&n&t&u&(1)&(2)&(4)&(5)\\\hline
13&0&1&1&1&0&1&0\\
37&1&1&2&2&1&3&0\\
61&2&3&2&4&1&3&2\\
109&4&3&3&6&3&6&3\\
157&6&3&4&8&5&9&4\\
181&7&5&3&10&5&8&7\\
229&9&5&3&12&7&10&9
\end{array}
\]

Dritter Fall; $p=24n+5$.
\[
\begin{aligned}
2t&\quad\text{Anzahl der Classen für den Determinans }-p;\\
2u&\quad\text{Anzahl der Classen für den Determinans }-3p.
\end{aligned}
\]

\SourceLeaf{301}
\[
\begin{aligned}
(1)=(2)=(6)=(7)=(11)=(12)&=n\\
(3)=(10)&=\tfrac12(2n+1+t)\\
(4)=(9)&=\tfrac12(2n-t+u)\\
(5)=(8)&=\tfrac12(2n+1-u)
\end{aligned}
\]
\[
\begin{array}{r|r|r|r|r|r|r|r}
p&n&t&u&(1)&(3)&(4)&(5)\\\hline
5&0&1&1&0&1&0&0\\
29&1&3&3&1&3&1&0\\
53&2&3&5&2&4&3&0\\
101&4&7&5&4&8&3&2\\
149&6&7&7&6&10&6&3\\
173&7&7&9&7&11&8&3\\
197&8&5&11&8&11&11&3\\
269&11&11&7&11&17&9&8
\end{array}
\]

Vierter Fall; $p=24n+17$.
\[
\begin{aligned}
2t&\quad\text{Anzahl der Classen für den Determinans }-p;\\
2u&\quad\text{Anzahl der Classen für den Determinans }-3p.
\end{aligned}
\]
\[
\begin{aligned}
(1)=(2)=(6)=(7)=(11)=(12)&=\tfrac16(6n+3+u)\\
(3)=(10)&=\tfrac16(6n+6+3t-2u)\\
(4)=(9)&=\tfrac16(6n+3-3t+u)\\
(5)=(8)&=\tfrac16(6n+6-2u)
\end{aligned}
\]
\[
\begin{array}{r|r|r|r|r|r|r|r}
p&n&t&u&(1)&(3)&(4)&(5)\\\hline
17&0&2&3&1&1&0&0\\
41&1&4&3&2&3&0&1\\
89&3&6&3&4&6&1&3\\
113&4&4&9&6&4&4&2\\
137&5&4&9&7&5&5&3\\
233&9&6&15&12&8&9&5\\
257&10&8&9&12&12&8&8
\end{array}
\]
\BookRule{.30\textwidth}
